% =====================================================================
%  Harald Høffding, review of Viktor Rydberg: Ny Testamentes Lærdomme om Forudtilværelsen og de sidste Ting.  Paa Dansk
%  ved Cand. theol. J. P. V. Petersen.  Kjøbenhavn 1881 (Andr. Schous Forlag).
%  In: Ude og Hjemme.  Nordisk illustreret Ugeblad, udgivet af Otto Borchsenius og F. Hendriksen.  Femte Aargang 1881-1882
%  (Kjøbenhavn: F. Hendriksen), Nr. 224, Søndagen den 15. Januar 1882, pp. 196-197.  Antiqua, two columns.
%
%  TRANSCRIPTION -- GENERATED; DO NOT EDIT BY HAND.  Built by ebook-method/rydberg-anmeldelse-1882/ (HOEFFDING-QUEUE #19):
%    sh ebook-method/rydberg-anmeldelse-1882/fetch.sh; python3 .../readings.py; python3 .../build.py   (from the repo root)
%
%  STATUS: COMPLETE (2026-10-17): the volume's title page and the whole review, pp. 196-197.
%
%  EDITION.  The only printing.  Text of record: Internet Archive ude-og-hjemme-aarg-5 (a private scan, 300 ppi JP2; the
%  same scan as the Böhme review, texts/hoeffding/boehme-anmeldelse-1882/), leaves 207-208; extract ~/bibliotek/Høffding,
%  Harald/1882-ude-og-hjemme-5-nr224-pp196-197.pdf (SCANS.tsv).  The volume's Indholdsfortegnelse: „Høffding, Harald: Viktor
%  Rydbergs »Ny Testamentes Lærdomme« .... 196.“  The review's head runs across both columns of p. 196 below the end of a
%  story („En Arrestation i Landemærket“); the review ends on p. 197, right column, above „To Billeder fra Grækenland“.
%
%  ACCURACY (2026-10-17; 1,700 words).
%  (a) No second copy found (as for the Böhme review: Internet Archive, Nasjonalbiblioteket, Runeberg).  The DRAFT is our
%      own tesseract reading (tessdata_best dan, 300 ppi, region by region: readings.py); the WITNESS is the Archive's own
%      tesseract hOCR, aligned line by line (70 of 171 lines differ, most of them alignment noise).
%  (b) EVERY LINE READ AT THE IMAGE in the main session against the draft (sheets.py: 171 line crops on eight sheets;
%      two 3x zooms): 9 lines corrected (reads.tsv): Cand. theol., Boström, lära, the note mark *), a speck over blive,
%      tørste (a turned e), saa, the note's (!) twice.  Head and signature in a bold grotesque; no italics, no
%      letterspacing in the review.
%  NOT done: a blind reading by a second reader (Hans asked for a lean session; the text is short).
%  PRINTED AS IS (misprints kept, marked at the line): „Fortolkning,.“ (p. 196), „Udganspunkt“ (p. 197).
%  CONVENTIONS.  Page turns: % --- p. N --- and \opage{N}; the columns are marked % (p. N, column c); the print's line
%  breaks are kept, a word divided at a line end is joined on the first line.  The note at its mark (the print's *)).
%  »...« as printed; — as ---.  The running heads are not transcribed.
% =====================================================================
\documentclass[12pt,a4paper,openany]{book}
\usepackage[utf8]{inputenc}
\usepackage[T1]{fontenc}
\usepackage{libertinus}
\usepackage{libertinust1math}
\usepackage[danish]{babel}
\usepackage{textalpha}   % Greek typed directly
\usepackage[protrusion=true,expansion=true]{microtype}
\usepackage{setspace}
\usepackage[margin=1.3in]{geometry}
\usepackage{fancyhdr}
\usepackage{hyperref}
\hypersetup{pdftitle={Anmeldelse af Viktor Rydberg, Ny Testamentes Lærdomme}, pdfauthor={Harald Høffding}, hidelinks}
\pagestyle{fancy}
\fancyhf{}
\fancyhead[LE,RO]{\thepage}
\fancyhead[C]{\textit{H. Høffding: Anmeldelse af Viktor Rydberg, Ny Testamentes Lærdomme (1882)}}
\fancypagestyle{plain}{\fancyhf{}\fancyhead[LE,RO]{\thepage}}
\setstretch{1.3}
\usepackage{marginnote}
\usepackage{graphicx}
\DeclareUnicodeCharacter{0254}{\reflectbox{c}}   % ɔ (id est)
\renewcommand{\thefootnote}{*}   % the print marks its one note *)
\makeatletter\renewcommand{\@makefnmark}{\hbox{\textsuperscript{\@thefnmark)}}}\makeatother
\newcommand{\opage}[1]{\marginnote{\footnotesize\textit{[#1]}}}
\newcommand{\signature}[1]{\par\smallskip\hfill #1\hspace*{2em}\par}
\pdfstringdefDisableCommands{\renewcommand{\footnote}[2][]{}\renewcommand{\opage}[1]{}}

\begin{document}

\begin{titlepage}
\begin{center}
% [a ruled frame round the page]
{\Huge UDE OG HJEMME.}\\[0.5em]
{\large NORDISK ILLUSTRERET UGEBLAD.}\\[1em]
% [an engraved cherub with a garland]
{\scriptsize UDGIVET AF}\\[0.3em]
{\sffamily OTTO BORCHSENIUS {\small OG} F. HENDRIKSEN.}\\[0.6em]
% [wavy rule]
{\small\bfseries FEMTE AARGANG 1881---1882.}\\
% [wavy rule]
\vfill
{\small KJØBENHAVN.}\\[0.3em]
{\scriptsize TRYKT PAA F. HENDRIKSENS FORLAG HOS J. JØRGENSEN \& CO.}
\end{center}
\end{titlepage}
\clearpage

% --- p. 196 ---
\setcounter{footnote}{0}%
\opage{196}\begin{center}{\sffamily\large Viktor Rydberg: Ny Testamentes Lærdomme om Forudtilværelsen og de sidste Ting.}\\[0.3em]{\small Paa Dansk ved Cand. theol. \textbf{\sffamily J. P. V. Petersen.} Kjøbenhavn 1881.  (Andr. Schous Forlag.)}\end{center}
% (p. 196, column 1)

Medens Viktor Rydberg er vel kjendt i store Kredse
som Digter, Oversætter og Kunsthistoriker, er det Færre
herhjemme, der vide, at han ogsaa indtager en fremragende
Plads som Religionshistoriker og hører til de meget faa
Forskere i de nordiske Lande, der, uden at høre hjemme
i den kirkelige Orthodoxis Rækker, med Iver og Begejstring
have hengivet sig til Studiet af de bibelske Skrifter. I Stedet
for bevidst eller ubevidst at lægge den gjældende Kirkelære
til Grund ved Fortolkningen af de nytestamentlige Bøger,
fordrer Rydberg, at man skal se dem i Lyset af den Tid,
i hvilken de bleve til, og lade de Tanker, der indeholdes
i dem, udfolde sig i fuld Frihed, idet man sér bort fra
den Maade, hvorpaa senere Tider have lagt sig dem tilrette.
% PRINTED AS IS (Fortolkning,.: p. 196: a comma and a stop, both printed (zoom1.png))
Det er med andre Ord en rent historisk Fortolkning,.
Rydberg tager Ordet for. Men saa meget han end i sin
Forsken holder sine egne Meninger og Sympathier udenfor,
er den blot historiske Forstaaelse dog ikke hans sidste
Maal. Det er hans Overbevisning, at jo mere man trænger
ned til Kildespringet for den store religiøse Bevægelse,
hvoraf hine gamle Skrifter opstod, des mere vil man ogsaa
finde Midler til at befri den religiøse Følelse fra den Strid
og Tvivl, der har udbredt sig i den aandelige Verden som
Følge af Sammenstødet mellem den kirkelige Theologi og
den videnskabelige Tænkning. Det Kristusbillede, de nytestamentlige
Skrifter give os, naar de opfattes i deres
eget Lys, indeholder efter hans Overbevisning ingen af de
Modsigelser, hvormed Kirkelæren er behæftet, men formaar
at tilfredsstille baade Tanken og Hjærtet. Ved at skyde
den kirkelige Orthodoxi til Side for den oprindelige nytestamentlige
Lære, som ogsaa herskede i Kirken, før
Kristendommen blev Statsreligion, vil Forfatteren give sin
Skjærv til Arbejdet paa at fortsætte Reformationen og
bygge »en Kirke, der opmuntrer Fornuften til at forske,
lige saavel som Hjærtet til at elske, idet den vil virkeliggjøre
% (p. 196, column 2)
baade Sandheden og Kjærligheden og gjennem begge
Friheden«. Forarbejder til en saadan Fortsættelse af, hvad
der indlededes ved Luthers Reformation, men altfor tidlig
standsedes, finder Forfatteren rundt om i den kristelige
Verden. I hans eget Fædreland have især »den fromme
og uforsagte Ignell og den store Tænker Boström« virket
i denne Retning.

Det foreliggende Skrift, som slutter sig til Forf.s større
Arbejde »Bibelns lära om Kristus«, der 1880 udkom i
fjerde Oplag, søger at godtgjøre, at Troen paa Menneskets
Forudtilværelse, som den gamle Verden havde ikke mindre
Trang til at antage end en Tilværelse efter Døden, og
som særlig gjennemtrængte hele den Tids religiøse Forestillingskreds,
paa hvilken de nytestamentlige Skrifter bleve
til, --- at denne Tro ogsaa er det nye Testamentes Lære.
Forf. støtter sig her især til Johannesevangeliet, de paulinske
Breve og Hebræerbrevet. De menneskelige Personligheder,
der i Sanseverdenen existere i Sondring og Ufuldkommenhed,
existerede i oprindelig Enhed i det evige Ideal, som
Johannesevangeliet kalder Logos og Paulus den himmelske
Adam. Ifølge Messiaslæren er det den historiske Kristus'
Opgave at forene det Sondrede og samle det Spredte.
Sanseverdenen er en Skole, gjennem hvilken den oprindelige,
ideale Enhed igjen skal virkeliggjøres. Den historiske
Udvikling afsluttes ifølge det nye Testamente med en
Verdensfornyelse, hvor hver enkelt Personlighed indtager
sin Plads i Aandernes ideale Verden. Læren om en evig
Fordømmelse, der har været herskende i Kirken siden det
femte Aarhundrede, erklærer Forf. for ubibelsk. Han nægter
ikke, at »Gjengjældelseslæren, hvorpaa den mosaiske Lov
hviler, som selv de største af Israels Profeter ikke fuldt
kunde overvinde, ogsaa kaster sin Skygge ind i det nye
Testament«; men han finder den stridende mod den
herskende Tankegang i de nytestamentlige Skrifter. Særlig
% --- p. 197 ---
\setcounter{footnote}{0}%
% (p. 197, column 1)
\opage{197}interessant er hans Paavisning af de Misforstaaelser, der
ere fremkomne ved urigtig Oversættelse og Fortolkning af
bibelske Udtryk (som naar »Hades« uden videre oversættes
ved »Helvede«, »aionios« ved »evig«). --- En nærmere Drøftelse
af disse Spørgsmaal skal jeg her ikke gaa ind paa.
En saadan vilde forudsætte en speciel exegetisk Lærdom,
som jeg ikke besidder. Jeg kan kun sige saa meget, at
Forf.s Fremstilling og Fortolkning synes utvungen og fører
Læseren ind i en Forestillingskreds, der synes at stemme
med den Tid og de Forhold, under hvilke de nytestamentlige
Skrifter bleve til. Forf.s historiske Sans og store
Sandhedskjærlighed\footnote[1]{En theologisk Modstander i Sverig har ikke kunnet gaa videre end at indrømme, at Rydberg »med en vis (!) Oprigtighed synes (!) at søge Sandheden«.} forbyder ham at raste, før han har
naaet de primitive Lag i den Jordbund, han undersøger.

Man kan ikke læse en Undersøgelse som den foreliggende
uden at tænke paa, hvorledes det vil gaa, dersom
de vordende Gejstliges Uddannelse skulde blive endnu
mindre videnskabelig og mere seminarieagtig end den for
Øjeblikket er. Dersom Rydberg har Ret, har falsk Oversættelse
af bibelske Udtryk ført til ubibelske Lærdomme i
Kirken. Men den ikke sproglig og kritisk dannede Theolog
maa holde sig til Oversættelsen og er udelukket fra at
korrigere slige Urigtigheder, ja endog fra at bedømme et
Indlæg som Rydbergs. Afgjørelsen af, om Rydberg har
Ret, beror ej paa, om han er »troende« eller ikke, men
om han har Ret i sin sproglige og historiske Fortolkning. ---

Viktor Rydbergs eget, religiøse Standpunkt falder, som
man sér af flere Udtalelser i hans Skrifter, kun for saa
vidt sammen med det, han betragter som det ægte bibelske,
som de bibelske Forestillinger om Forudtilværelsen, timelig
og sanselig Tilværelse og Verdensfornyelse kunne tjene til
Symboler for religiøse Ideer, der i sig selv pege ud over
Tidens og Rummets Former. Rydberg slutter sig nemlig
til den i Sverig herskende idealistiske Filosofi, efter hvis
Lære det Guddommelige ej kan være Vorden og Udvikling
underlagt. Kristendommens Opfattelse af Verdenslivet som
en stor historisk Udviklingsproces, i hvilken Guddommen
selv virker og lider med, kan for et saadant Standpunkt
kun have symbolsk Værdi. Selv Hegel er jo de svenske
Idealister for empirisk. Her forekommer det mig, at man
af den svenske Lessing kunde have ventet et friere Blik.
Jeg bebrejder ham ikke, at han i de bibelske Forestillinger
sér en poetisk Symbolik og ikke kan gaa ind paa den
bogstavelige Opfattelse af dem. Saaledes stiller jeg mig
ogsaa selv til disse ærværdige Overleveringer, idet jeg forresten
finder Summen af den bibelske Visdom, ikke i nogen
som helst af de dogmatiske Forestillinger, men i Forkyndelsen
af den ubetingede Menneskekjærligheds Evangelium
og i Saligprisningen af dem, der hungre og tørste
efter Retfærdighed. Men jeg tror, at Rydberg kunde være
gaaet et Skridt videre og have anerkjendt det som betydningsfuldt,
at de kristelige Grundskrifter paa deres Vis anerkjende
Udviklingen som en Verdenslov. Naar den moderne
Videnskab fra saa mange Sider kommer til den samme
Tanke, vilde her jo for den, der vil opsøge det Fælles bag
% PRINTED AS IS (Udganspunkt: p. 197: for Udgangspunkt (ryd-L-06))
de stridende Modsætninger, være et frugtbart Udganspunkt.
En Idealisme, der betragter Tid og Udvikling som blotte
Illusioner, spærrer sig ude fra alt frugtbart Samkvem med
det virkelige Livs Rørelser. Det er muligt, at Begreber
som Tid og Udvikling, lige saa vel som alle andre Begreber,
kun have Gyldighed fra et menneskeligt Standpunkt,
men paa et andet kunne vi nu en Gang ikke stille os.
Den kritiske Indsigt i vore Grundbegrebers Sammenhæng
med vor Organisation og vor Plads i Verden er ikke det
Samme som en Overbevisning om, at de ere illusoriske;
tværtimod: hænge vore Grundbegreber virkelig sammen
med vor Natur, saa have de ogsaa al den Gyldighed, vi i
det Hele kunne tillægge Begreber. --- Dersom Rydberg
ikke gik ud fra en saa abstrakt Idealisme, vilde han ogsaa
med større Sympathi følge de nyere Bestræbelser paa det
psykologiske og filosofiske Omraade. Der er her aabenbart
en Mangel i hans rige og mangesidige Dannelse. Medens
% (p. 197, column 2)
hans Tanke er sleben til Forstaaelse og Behandling af
digteriske, religiøse og spekulative Forestillingsformer, har
han næppe med samme Ihærdighed og Sympathi sat sig
ind i Erfaringsvidenskabernes Methode og dens Betydning
ogsaa for Filosofien. Hvor fremmed han er for den rent
erfaringsmæssige (induktive) Synsmaade, faar man et Indtryk
af, naar han lejlighedsvis omtaler den yngre Fichte og
Carus som Forskere, der »uafhængig af metafysiske Theorier
have forsøgt ad naturvidenskabelig Vej at gjennemtrænge
Antropologiens Problem«. De nævnte Mænd beraabe sig
ganske vist meget ofte paa »Erfaring«; men dertil regne
de da allehaande uklar Mystik og spiritistisk Fantasteri,
naar det passer efter deres Meninger. Den virkelige
Erfaringsmethode laa dem fjærnt. (Smlgn. Lotze's Stridsskrift
mod den yngre Fichte).

Disse Spørgsmaal berøres kun rent antydningsvis i
Bogen, skjønt Antydningerne gjøre den filosofiske Læser
nysgjerrig. Det vilde sikkert være af stor Betydning for
dem i Norden, der hver paa sin Maade hylde en friere
religiøs Overbevisning, om en Mand som Viktor Rydberg
vilde give os en selvstændig Behandling af det religiøse
Problem. Med sin klare Tanke, sit herlige Sprog, og
fremfor Alt sin varme Følelse for alt Ædelt og Sandt har
han vigtige Betingelser for at løse en saadan Opgave.
Hans utrættelige Forsken vilde da vistnok ogsaa føre ham
til at finde adskilligt brugbart Malm ogsaa i dem af Nutidens
filosofiske Bestræbelser, i hvilke han, som det synes, hidtil
kun har sét Slakker.
\signature{\textsf{\bfseries Harald Høffding.}}

\end{document}
